\PassOptionsToPackage{dvipsnames,table}{xcolor}
\documentclass[11pt, letterpaper]{article}
\newcommand{\revision}[1]{#1}

% packages
\usepackage[left=1in, right=1in, top=1in, bottom=1in]{geometry}
\usepackage[utf8]{inputenc}
\usepackage[T1]{fontenc}
\usepackage{amsmath,amssymb,amsthm,amsfonts}
\usepackage{bm}
\usepackage{graphicx}
\usepackage[font=footnotesize,labelfont=bf,figurename=Figure,labelformat=empty,labelsep=none]{caption}
\usepackage{xcolor}
\usepackage{hyperref}
\usepackage{booktabs}
\usepackage[sort&compress,numbers]{natbib}
%%%%% NEW MATH DEFINITIONS %%%%%

% Mark sections of captions for referring to divisions of figures
\newcommand{\figleft}{{\em (Left)}}
\newcommand{\figcenter}{{\em (Center)}}
\newcommand{\figright}{{\em (Right)}}
\newcommand{\figtop}{{\em (Top)}}
\newcommand{\figbottom}{{\em (Bottom)}}
\newcommand{\captiona}{{\em (a)}}
\newcommand{\captionb}{{\em (b)}}
\newcommand{\captionc}{{\em (c)}}
\newcommand{\captiond}{{\em (d)}}
\newcommand{\rev}[1]{\textcolor{black}{#1}}% 红色
% Highlight a newly defined term
\newcommand{\newterm}[1]{{\bf #1}}


% Figure reference, lower-case.
\def\figref#1{figure~\ref{#1}}
% Figure reference, capital. For start of sentence
\def\Figref#1{Figure~\ref{#1}}
\def\twofigref#1#2{figures \ref{#1} and \ref{#2}}
\def\quadfigref#1#2#3#4{figures \ref{#1}, \ref{#2}, \ref{#3} and \ref{#4}}
% Section reference, lower-case.
\def\secref#1{section~\ref{#1}}
% Section reference, capital.
\def\Secref#1{Section~\ref{#1}}
% Reference to two sections.
\def\twosecrefs#1#2{sections \ref{#1} and \ref{#2}}
% Reference to three sections.
\def\secrefs#1#2#3{sections \ref{#1}, \ref{#2} and \ref{#3}}
% Reference to an equation, lower-case.
\def\eqref#1{equation~\ref{#1}}
% Reference to an equation, upper case
\def\Eqref#1{Equation~\ref{#1}}
% \def\Eqref#1{Eq.~\ref{#1}
% A raw reference to an equation---avoid using if possible
\def\plaineqref#1{\ref{#1}}
% Reference to a chapter, lower-case.
\def\chapref#1{chapter~\ref{#1}}
% Reference to an equation, upper case.
\def\Chapref#1{Chapter~\ref{#1}}
% Reference to a range of chapters
\def\rangechapref#1#2{chapters\ref{#1}--\ref{#2}}
% Reference to an algorithm, lower-case.
\def\algref#1{algorithm~\ref{#1}}
% Reference to an algorithm, upper case.
\def\Algref#1{Algorithm~\ref{#1}}
\def\twoalgref#1#2{algorithms \ref{#1} and \ref{#2}}
\def\Twoalgref#1#2{Algorithms \ref{#1} and \ref{#2}}
% Reference to a part, lower case
\def\partref#1{part~\ref{#1}}
% Reference to a part, upper case
\def\Partref#1{Part~\ref{#1}}
\def\twopartref#1#2{parts \ref{#1} and \ref{#2}}

\def\ceil#1{\lceil #1 \rceil}
\def\floor#1{\lfloor #1 \rfloor}
\def\1{\bm{1}}
\newcommand{\train}{\mathcal{D}}
\newcommand{\valid}{\mathcal{D_{\mathrm{valid}}}}
\newcommand{\test}{\mathcal{D_{\mathrm{test}}}}

\def\eps{{\epsilon}}


% Random variables
\def\reta{{\textnormal{$\eta$}}}
\def\ra{{\textnormal{a}}}
\def\rb{{\textnormal{b}}}
\def\rc{{\textnormal{c}}}
\def\rd{{\textnormal{d}}}
\def\re{{\textnormal{e}}}
\def\rf{{\textnormal{f}}}
\def\rg{{\textnormal{g}}}
\def\rh{{\textnormal{h}}}
\def\ri{{\textnormal{i}}}
\def\rj{{\textnormal{j}}}
\def\rk{{\textnormal{k}}}
\def\rl{{\textnormal{l}}}
% rm is already a command, just don't name any random variables m
\def\rn{{\textnormal{n}}}
\def\ro{{\textnormal{o}}}
\def\rp{{\textnormal{p}}}
\def\rq{{\textnormal{q}}}
\def\rr{{\textnormal{r}}}
\def\rs{{\textnormal{s}}}
\def\rt{{\textnormal{t}}}
\def\ru{{\textnormal{u}}}
\def\rv{{\textnormal{v}}}
\def\rw{{\textnormal{w}}}
\def\rx{{\textnormal{x}}}
\def\ry{{\textnormal{y}}}
\def\rz{{\textnormal{z}}}

% Random vectors
\def\rvepsilon{{\mathbf{\epsilon}}}
\def\rvtheta{{\mathbf{\theta}}}
\def\rva{{\mathbf{a}}}
\def\rvb{{\mathbf{b}}}
\def\rvc{{\mathbf{c}}}
\def\rvd{{\mathbf{d}}}
\def\rve{{\mathbf{e}}}
\def\rvf{{\mathbf{f}}}
\def\rvg{{\mathbf{g}}}
\def\rvh{{\mathbf{h}}}
\def\rvu{{\mathbf{i}}}
\def\rvj{{\mathbf{j}}}
\def\rvk{{\mathbf{k}}}
\def\rvl{{\mathbf{l}}}
\def\rvm{{\mathbf{m}}}
\def\rvn{{\mathbf{n}}}
\def\rvo{{\mathbf{o}}}
\def\rvp{{\mathbf{p}}}
\def\rvq{{\mathbf{q}}}
\def\rvr{{\mathbf{r}}}
\def\rvs{{\mathbf{s}}}
\def\rvt{{\mathbf{t}}}
\def\rvu{{\mathbf{u}}}
\def\rvv{{\mathbf{v}}}
\def\rvw{{\mathbf{w}}}
\def\rvx{{\mathbf{x}}}
\def\rvy{{\mathbf{y}}}
\def\rvz{{\mathbf{z}}}

% Elements of random vectors
\def\erva{{\textnormal{a}}}
\def\ervb{{\textnormal{b}}}
\def\ervc{{\textnormal{c}}}
\def\ervd{{\textnormal{d}}}
\def\erve{{\textnormal{e}}}
\def\ervf{{\textnormal{f}}}
\def\ervg{{\textnormal{g}}}
\def\ervh{{\textnormal{h}}}
\def\ervi{{\textnormal{i}}}
\def\ervj{{\textnormal{j}}}
\def\ervk{{\textnormal{k}}}
\def\ervl{{\textnormal{l}}}
\def\ervm{{\textnormal{m}}}
\def\ervn{{\textnormal{n}}}
\def\ervo{{\textnormal{o}}}
\def\ervp{{\textnormal{p}}}
\def\ervq{{\textnormal{q}}}
\def\ervr{{\textnormal{r}}}
\def\ervs{{\textnormal{s}}}
\def\ervt{{\textnormal{t}}}
\def\ervu{{\textnormal{u}}}
\def\ervv{{\textnormal{v}}}
\def\ervw{{\textnormal{w}}}
\def\ervx{{\textnormal{x}}}
\def\ervy{{\textnormal{y}}}
\def\ervz{{\textnormal{z}}}

% Random matrices
\def\rmA{{\mathbf{A}}}
\def\rmB{{\mathbf{B}}}
\def\rmC{{\mathbf{C}}}
\def\rmD{{\mathbf{D}}}
\def\rmE{{\mathbf{E}}}
\def\rmF{{\mathbf{F}}}
\def\rmG{{\mathbf{G}}}
\def\rmH{{\mathbf{H}}}
\def\rmI{{\mathbf{I}}}
\def\rmJ{{\mathbf{J}}}
\def\rmK{{\mathbf{K}}}
\def\rmL{{\mathbf{L}}}
\def\rmM{{\mathbf{M}}}
\def\rmN{{\mathbf{N}}}
\def\rmO{{\mathbf{O}}}
\def\rmP{{\mathbf{P}}}
\def\rmQ{{\mathbf{Q}}}
\def\rmR{{\mathbf{R}}}
\def\rmS{{\mathbf{S}}}
\def\rmT{{\mathbf{T}}}
\def\rmU{{\mathbf{U}}}
\def\rmV{{\mathbf{V}}}
\def\rmW{{\mathbf{W}}}
\def\rmX{{\mathbf{X}}}
\def\rmY{{\mathbf{Y}}}
\def\rmZ{{\mathbf{Z}}}

% Elements of random matrices
\def\ermA{{\textnormal{A}}}
\def\ermB{{\textnormal{B}}}
\def\ermC{{\textnormal{C}}}
\def\ermD{{\textnormal{D}}}
\def\ermE{{\textnormal{E}}}
\def\ermF{{\textnormal{F}}}
\def\ermG{{\textnormal{G}}}
\def\ermH{{\textnormal{H}}}
\def\ermI{{\textnormal{I}}}
\def\ermJ{{\textnormal{J}}}
\def\ermK{{\textnormal{K}}}
\def\ermL{{\textnormal{L}}}
\def\ermM{{\textnormal{M}}}
\def\ermN{{\textnormal{N}}}
\def\ermO{{\textnormal{O}}}
\def\ermP{{\textnormal{P}}}
\def\ermQ{{\textnormal{Q}}}
\def\ermR{{\textnormal{R}}}
\def\ermS{{\textnormal{S}}}
\def\ermT{{\textnormal{T}}}
\def\ermU{{\textnormal{U}}}
\def\ermV{{\textnormal{V}}}
\def\ermW{{\textnormal{W}}}
\def\ermX{{\textnormal{X}}}
\def\ermY{{\textnormal{Y}}}
\def\ermZ{{\textnormal{Z}}}

% Vectors
\def\vzero{{\bm{0}}}
\def\vone{{\bm{1}}}
\def\vmu{{\bm{\mu}}}
\def\vtheta{{\bm{\theta}}}
\def\va{{\bm{a}}}
\def\vb{{\bm{b}}}
\def\vc{{\bm{c}}}
\def\vd{{\bm{d}}}
\def\ve{{\bm{e}}}
\def\vf{{\bm{f}}}
\def\vg{{\bm{g}}}
\def\vh{{\bm{h}}}
\def\vi{{\bm{i}}}
\def\vj{{\bm{j}}}
\def\vk{{\bm{k}}}
\def\vl{{\bm{l}}}
\def\vm{{\bm{m}}}
\def\vn{{\bm{n}}}
\def\vo{{\bm{o}}}
\def\vp{{\bm{p}}}
\def\vq{{\bm{q}}}
\def\vr{{\bm{r}}}
\def\vs{{\bm{s}}}
\def\vt{{\bm{t}}}
\def\vu{{\bm{u}}}
\def\vv{{\bm{v}}}
\def\vw{{\bm{w}}}
\def\vx{{\bm{x}}}
\def\vy{{\bm{y}}}
\def\vz{{\bm{z}}}
\def\vepsilon{{\bm{\epsilon}}}

% Elements of vectors
\def\evalpha{{\alpha}}
\def\evbeta{{\beta}}
\def\evepsilon{{\epsilon}}
\def\evlambda{{\lambda}}
\def\evomega{{\omega}}
\def\evmu{{\mu}}
\def\evpsi{{\psi}}
\def\evsigma{{\sigma}}
\def\evtheta{{\theta}}
\def\eva{{a}}
\def\evb{{b}}
\def\evc{{c}}
\def\evd{{d}}
\def\eve{{e}}
\def\evf{{f}}
\def\evg{{g}}
\def\evh{{h}}
\def\evi{{i}}
\def\evj{{j}}
\def\evk{{k}}
\def\evl{{l}}
\def\evm{{m}}
\def\evn{{n}}
\def\evo{{o}}
\def\evp{{p}}
\def\evq{{q}}
\def\evr{{r}}
\def\evs{{s}}
\def\evt{{t}}
\def\evu{{u}}
\def\evv{{v}}
\def\evw{{w}}
\def\evx{{x}}
\def\evy{{y}}
\def\evz{{z}}

% Matrix
\def\mA{{\bm{A}}}
\def\mB{{\bm{B}}}
\def\mC{{\bm{C}}}
\def\mD{{\bm{D}}}
\def\mE{{\bm{E}}}
\def\mF{{\bm{F}}}
\def\mG{{\bm{G}}}
\def\mH{{\bm{H}}}
\def\mI{{\bm{I}}}
\def\mJ{{\bm{J}}}
\def\mK{{\bm{K}}}
\def\mL{{\bm{L}}}
\def\mM{{\bm{M}}}
\def\mN{{\bm{N}}}
\def\mO{{\bm{O}}}
\def\mP{{\bm{P}}}
\def\mQ{{\bm{Q}}}
\def\mR{{\bm{R}}}
\def\mS{{\bm{S}}}
\def\mT{{\bm{T}}}
\def\mU{{\bm{U}}}
\def\mV{{\bm{V}}}
\def\mW{{\bm{W}}}
\def\mX{{\bm{X}}}
\def\mY{{\bm{Y}}}
\def\mZ{{\bm{Z}}}
\def\mBeta{{\bm{\beta}}}
\def\mPhi{{\bm{\Phi}}}
\def\mLambda{{\bm{\Lambda}}}
\def\mSigma{{\bm{\Sigma}}}

% Tensor
\DeclareMathAlphabet{\mathsfit}{\encodingdefault}{\sfdefault}{m}{sl}
\SetMathAlphabet{\mathsfit}{bold}{\encodingdefault}{\sfdefault}{bx}{n}
\newcommand{\tens}[1]{\bm{\mathsfit{#1}}}
\def\tA{{\tens{A}}}
\def\tB{{\tens{B}}}
\def\tC{{\tens{C}}}
\def\tD{{\tens{D}}}
\def\tE{{\tens{E}}}
\def\tF{{\tens{F}}}
\def\tG{{\tens{G}}}
\def\tH{{\tens{H}}}
\def\tI{{\tens{I}}}
\def\tJ{{\tens{J}}}
\def\tK{{\tens{K}}}
\def\tL{{\tens{L}}}
\def\tM{{\tens{M}}}
\def\tN{{\tens{N}}}
\def\tO{{\tens{O}}}
\def\tP{{\tens{P}}}
\def\tQ{{\tens{Q}}}
\def\tR{{\tens{R}}}
\def\tS{{\tens{S}}}
\def\tT{{\tens{T}}}
\def\tU{{\tens{U}}}
\def\tV{{\tens{V}}}
\def\tW{{\tens{W}}}
\def\tX{{\tens{X}}}
\def\tY{{\tens{Y}}}
\def\tZ{{\tens{Z}}}


% Graph
\def\gA{{\mathcal{A}}}
\def\gB{{\mathcal{B}}}
\def\gC{{\mathcal{C}}}
\def\gD{{\mathcal{D}}}
\def\gE{{\mathcal{E}}}
\def\gF{{\mathcal{F}}}
\def\gG{{\mathcal{G}}}
\def\gH{{\mathcal{H}}}
\def\gI{{\mathcal{I}}}
\def\gJ{{\mathcal{J}}}
\def\gK{{\mathcal{K}}}
\def\gL{{\mathcal{L}}}
\def\gM{{\mathcal{M}}}
\def\gN{{\mathcal{N}}}
\def\gO{{\mathcal{O}}}
\def\gP{{\mathcal{P}}}
\def\gQ{{\mathcal{Q}}}
\def\gR{{\mathcal{R}}}
\def\gS{{\mathcal{S}}}
\def\gT{{\mathcal{T}}}
\def\gU{{\mathcal{U}}}
\def\gV{{\mathcal{V}}}
\def\gW{{\mathcal{W}}}
\def\gX{{\mathcal{X}}}
\def\gY{{\mathcal{Y}}}
\def\gZ{{\mathcal{Z}}}

% Sets
\def\sA{{\mathbb{A}}}
\def\sB{{\mathbb{B}}}
\def\sC{{\mathbb{C}}}
\def\sD{{\mathbb{D}}}
% Don't use a set called E, because this would be the same as our symbol
% for expectation.
\def\sF{{\mathbb{F}}}
\def\sG{{\mathbb{G}}}
\def\sH{{\mathbb{H}}}
\def\sI{{\mathbb{I}}}
\def\sJ{{\mathbb{J}}}
\def\sK{{\mathbb{K}}}
\def\sL{{\mathbb{L}}}
\def\sM{{\mathbb{M}}}
\def\sN{{\mathbb{N}}}
\def\sO{{\mathbb{O}}}
\def\sP{{\mathbb{P}}}
\def\sQ{{\mathbb{Q}}}
\def\sR{{\mathbb{R}}}
\def\sS{{\mathbb{S}}}
\def\sT{{\mathbb{T}}}
\def\sU{{\mathbb{U}}}
\def\sV{{\mathbb{V}}}
\def\sW{{\mathbb{W}}}
\def\sX{{\mathbb{X}}}
\def\sY{{\mathbb{Y}}}
\def\sZ{{\mathbb{Z}}}

% Entries of a matrix
\def\emLambda{{\Lambda}}
\def\emA{{A}}
\def\emB{{B}}
\def\emC{{C}}
\def\emD{{D}}
\def\emE{{E}}
\def\emF{{F}}
\def\emG{{G}}
\def\emH{{H}}
\def\emI{{I}}
\def\emJ{{J}}
\def\emK{{K}}
\def\emL{{L}}
\def\emM{{M}}
\def\emN{{N}}
\def\emO{{O}}
\def\emP{{P}}
\def\emQ{{Q}}
\def\emR{{R}}
\def\emS{{S}}
\def\emT{{T}}
\def\emU{{U}}
\def\emV{{V}}
\def\emW{{W}}
\def\emX{{X}}
\def\emY{{Y}}
\def\emZ{{Z}}
\def\emSigma{{\Sigma}}

% entries of a tensor
% Same font as tensor, without \bm wrapper
\newcommand{\etens}[1]{\mathsfit{#1}}
\def\etLambda{{\etens{\Lambda}}}
\def\etA{{\etens{A}}}
\def\etB{{\etens{B}}}
\def\etC{{\etens{C}}}
\def\etD{{\etens{D}}}
\def\etE{{\etens{E}}}
\def\etF{{\etens{F}}}
\def\etG{{\etens{G}}}
\def\etH{{\etens{H}}}
\def\etI{{\etens{I}}}
\def\etJ{{\etens{J}}}
\def\etK{{\etens{K}}}
\def\etL{{\etens{L}}}
\def\etM{{\etens{M}}}
\def\etN{{\etens{N}}}
\def\etO{{\etens{O}}}
\def\etP{{\etens{P}}}
\def\etQ{{\etens{Q}}}
\def\etR{{\etens{R}}}
\def\etS{{\etens{S}}}
\def\etT{{\etens{T}}}
\def\etU{{\etens{U}}}
\def\etV{{\etens{V}}}
\def\etW{{\etens{W}}}
\def\etX{{\etens{X}}}
\def\etY{{\etens{Y}}}
\def\etZ{{\etens{Z}}}

% The true underlying data generating distribution
\newcommand{\pdata}{p_{\rm{data}}}
% The empirical distribution defined by the training set
\newcommand{\ptrain}{\hat{p}_{\rm{data}}}
\newcommand{\Ptrain}{\hat{P}_{\rm{data}}}
% The model distribution
\newcommand{\pmodel}{p_{\rm{model}}}
\newcommand{\Pmodel}{P_{\rm{model}}}
\newcommand{\ptildemodel}{\tilde{p}_{\rm{model}}}
% Stochastic autoencoder distributions
\newcommand{\pencode}{p_{\rm{encoder}}}
\newcommand{\pdecode}{p_{\rm{decoder}}}
\newcommand{\precons}{p_{\rm{reconstruct}}}

\newcommand{\laplace}{\mathrm{Laplace}} % Laplace distribution

\newcommand{\E}{\mathbb{E}}
\newcommand{\Ls}{\mathcal{L}}
\newcommand{\R}{\mathbb{R}}
\newcommand{\emp}{\tilde{p}}
\newcommand{\lr}{\alpha}
\newcommand{\reg}{\lambda}
\newcommand{\rect}{\mathrm{rectifier}}
\newcommand{\softmax}{\mathrm{softmax}}
\newcommand{\sigmoid}{\sigma}
\newcommand{\softplus}{\zeta}
\newcommand{\KL}{D_{\mathrm{KL}}}
\newcommand{\Var}{\mathrm{Var}}
\newcommand{\standarderror}{\mathrm{SE}}
\newcommand{\Cov}{\mathrm{Cov}}
% Wolfram Mathworld says $L^2$ is for function spaces and $\ell^2$ is for vectors
% But then they seem to use $L^2$ for vectors throughout the site, and so does
% wikipedia.
\newcommand{\normlzero}{L^0}
\newcommand{\normlone}{L^1}
\newcommand{\normltwo}{L^2}
\newcommand{\normlp}{L^p}
\newcommand{\normmax}{L^\infty}

\newcommand{\parents}{Pa} % See usage in notation.tex. Chosen to match Daphne's book.

\DeclareMathOperator*{\argmax}{arg\,max}
\DeclareMathOperator*{\argmin}{arg\,min}

\DeclareMathOperator{\sign}{sign}
\DeclareMathOperator{\Tr}{Tr}
\let\ab\allowbreak

\newcommand{\redz}[2]{\textcolor{black}{#1 [@zeng: #2]}}
\newcommand{\reddelz}[2]{\textcolor{black}{#1 [@zeng delete: #2]}}

\graphicspath{{./figures/}{./Figure/}}
\newcommand{\eref}[1]{(\ref{#1})}
\newcommand{\graph}[1]{\mathcal{}}
\providecommand{\keywords}[1]{\textbf{\textit{Keywords: }} #1}
\newcommand{\invdesroute}{InvDesRoute}

\hypersetup{
unicode=false,
pdftoolbar=true,
pdfmenubar=true,
pdffitwindow=false,
pdfstartview={FitH},
  pdftitle={InvDesRoute: route-network planning for inorganic synthesis},
pdfauthor={Yiheng Zhan, Ze-Feng Gao, Zhong-Yi Lu},
pdfsubject={InvDesRoute: route-network planning for inorganic synthesis},
pdfcreator={Creator},
pdfproducer={Producer},
pdfkeywords={inorganic synthesis, computational materials science, route planning, reaction networks, synthesizability},
pdfnewwindow=true,
colorlinks=true,
linkcolor=blue,
citecolor=blue,
filecolor=blue,
urlcolor=blue
}


\begin{document}

\title{\textbf{\invdesroute{}:\ route-network planning for inorganic synthesis}}

\author{Yi-Heng Zhan$^{1,\dagger}$, Ze-Feng Gao$^{1,\dagger,*}$, Zhong-Yi Lu$^{1,*}$\\
\small $^1$School of Physics and Key Laboratory of Quantum State Construction \\
\small and Manipulation (Ministry of Education), Renmin University of China, Beijing 100872, China\\
\small $^\dagger$These authors contributed equally.\\
\small $^*$e-mail: zfgao@ruc.edu.cn; zlu@ruc.edu.cn}
\date{}
\maketitle


\date{}

\maketitle

\normalsize

\vspace{-28pt} 
\newcommand{\mpform}{MP-form}
\begin{abstract}
\small

Most computational materials screens stop at whether a target crystal is stable, whereas experimental synthesis starts with which route to try, why that route is plausible and what uncertainty should be refined before effort is spent. Here we report \invdesroute, a route-network framework that turns formulas or user-supplied structures into ranked, checkable inorganic synthesis hypotheses. The framework links Materials Project phase context, precursor and gas reservoirs, PET-OAM-XL energies, balanced formation and decomposition reactions, route-family similarity and refinement flags. In a 120-target benchmark, \invdesroute{} enumerates 18,063 formation reactions and achieves 0.931 PET/MP-form sign consistency, with top-ranked route agreement for 99 targets. Against 145 DOI-backed text-mined literature routes, full networks contain moderate-or-better route-family analogues for 143 routes, whereas PET-energy-only ranking recovers only 13 in the top-5 window and \invdesroute{} scoring with diversity-aware display recovers 113. \invdesroute{} therefore makes inorganic synthesizability a route-network decision rather than a scalar target-stability label, with release records that separate present retrospective claims from future blind validation, calibration and reproduction.
% 中文：大多数计算材料筛选停留在判断目标晶体是否稳定，而实验合成首先要决定尝试哪条路线、为什么该路线合理，以及在投入工作之前应优先精修哪些不确定性。本文报道 \invdesroute，一个路线网络框架，可将化学式或用户提供的结构转化为经过排序、可检查的无机合成假设。该框架连接 Materials Project 相背景、前驱体和气体库、PET-OAM-XL 能量、配平的形成与分解反应、路线族相似性和精修标记。在 120 个目标的 benchmark 中，\invdesroute{} 枚举 18,063 条形成反应，并达到 0.931 的 PET/MP-form 符号一致性，其中 99 个目标的最高排序路线一致。与 145 条带 DOI 的文本挖掘文献路线相比，完整网络为 143 条路线包含 moderate-or-better 的路线族类似物，而仅按 PET 能量排序在 top-5 窗口中只恢复 13 条；采用 \invdesroute{} 评分并结合 diversity-aware 展示可恢复 113 条。因此，\invdesroute{} 将无机可合成性从标量目标稳定性标签转化为路线网络决策，并通过发布记录区分当前回顾性主张与未来的盲法验证、校准和复现。

\vspace{4pt}
\end{abstract}

\keywords{inorganic synthesis; computational materials science; route-network planning; reaction networks; synthesizability; machine-learning interatomic potentials}
% 中文：关键词：无机合成；计算材料科学；路线网络规划；反应网络；可合成性；机器学习原子间势

\section*{Introduction}
% 中文：引言

Computational materials discovery can now screen crystals across large chemical spaces using phase diagrams, formation energies, crystal databases, substitution rules and generative models \cite{hellenbrandt2004icsd,jain2013materialsproject,jain2016computational_predictions,hautier2011ionic_substitutions,merchant2023gnome,zeni2025mattergen}. Descriptor toolkits, crystal graph models and benchmark suites have broadened this ecosystem from hand-built representations to structure-aware property prediction with predefined tasks and test splits \cite{matminer2018,xie2018cgcnn,chen2019megnet,dunn2020matbench}. Stability- and synthesizability-focused benchmarks further show that formation-energy accuracy, stability classification and target-level synthesizability are related but distinct tasks \cite{bartel2020critical_stability,jang2020structure_synthesizability,davariashtiyani2021synthesizability,riebesell2025matbench_discovery}. Yet an experimental synthesis decision asks a different question: which reactants to try, which byproducts and reservoirs matter, which competing phases threaten the target and what uncertainty should be refined next.
% 中文：如今，计算材料发现已经能够利用相图、形成能、晶体数据库、取代规则和生成模型，在大规模化学空间中筛选晶体 \cite{hellenbrandt2004icsd,jain2013materialsproject,jain2016computational_predictions,hautier2011ionic_substitutions,merchant2023gnome,zeni2025mattergen}。描述符工具包、晶体图模型和 benchmark 套件进一步扩展了这一生态，使其从手工构建的表征发展到带有预定义任务和测试划分的结构感知性质预测 \cite{matminer2018,xie2018cgcnn,chen2019megnet,dunn2020matbench}。面向稳定性和可合成性的 benchmark 进一步表明，形成能精度、稳定性分类和目标级可合成性是相关但彼此不同的任务 \cite{bartel2020critical_stability,jang2020structure_synthesizability,davariashtiyani2021synthesizability,riebesell2025matbench_discovery}。然而，实验合成决策提出的是另一个问题：应该尝试哪些反应物，哪些副产物和库物种重要，哪些竞争相会威胁目标，以及下一步应精修哪类不确定性。

The missing object is the synthesis route network. Unlike molecular retrosynthesis, where neural and symbolic planners can operate on discrete molecular reactions \cite{segler2018synthesis_planning}, inorganic routes couple periodic solids, multicomponent phase competition, volatile reservoirs, precursor decomposition, metastability and condition-dependent thermochemistry \cite{sun2016metastability,aykol2018metastable_limit,mcdermott2023selectivity}. Reaction-network analyses and solid-state route planners have made this graph computational \cite{aykol2019network_analysis,aykol2021routes,mcdermott2021graph}. In parallel, text-mined corpora, literature-trained synthesis models, precursor and condition predictors, LLM or generative planners and autonomous laboratories show that route choices, experimental priors and negative evidence must be represented explicitly \cite{olivetti2020nlp_materials,kononova2019textmined,kim2020literature_synthesis_planning,huo2022conditions,he2023precursor,szymanski2023arrows3,kim2024llm,pan2026diffsyn,szymanski2023alab}. They also expose an evaluation problem: exact reaction matching is brittle because carbonates, hydrates, nitrates, hydroxides, oxalates or ammonium salts can express the same route family after dehydration, gas release or decomposition.
% 中文：缺失的对象是合成路线网络。不同于分子逆合成中神经规划器和符号规划器可以作用于离散分子反应 \cite{segler2018synthesis_planning}，无机路线同时耦合周期性固体、多组分相竞争、挥发性库物种、前驱体分解、亚稳性以及依赖条件的热化学 \cite{sun2016metastability,aykol2018metastable_limit,mcdermott2023selectivity}。反应网络分析和固相路线规划器已经使这种图结构可以被计算处理 \cite{aykol2019network_analysis,aykol2021routes,mcdermott2021graph}。与此同时，文本挖掘语料库、基于文献训练的合成模型、前驱体和条件预测器、LLM 或生成式规划器以及自主实验室都表明，路线选择、实验先验和负证据必须被显式表示 \cite{olivetti2020nlp_materials,kononova2019textmined,kim2020literature_synthesis_planning,huo2022conditions,he2023precursor,szymanski2023arrows3,kim2024llm,pan2026diffsyn,szymanski2023alab}。它们也暴露出一个评价问题：精确反应匹配很脆弱，因为碳酸盐、水合物、硝酸盐、氢氧化物、草酸盐或铵盐在脱水、放气或分解后可能表达同一类路线族。

We introduce \invdesroute{} to make this route network the primary computational unit. A target formula or user-supplied CIF is expanded into a phase context, enumerated formation and decomposition reactions, scalable PET-OAM-XL energies and precursor-aware recommendations with refinement flags. The contribution is a unit shift: from target-stability labels to checkable route hypotheses linking structure, phase space, route chemistry and uncertainty.
% 中文：我们提出 \invdesroute{}，将这种路线网络作为主要计算单元。一个目标化学式或用户提供的 CIF 会被扩展为相背景、枚举的形成与分解反应、可扩展的 PET-OAM-XL 能量，以及带有精修标记的前驱体感知推荐。本文的贡献是计算单元的转变：从目标稳定性标签转向可检查的路线假设，并将结构、相空间、路线化学和不确定性连接起来。

We evaluate that object through energy-engine tests, benchmark-scale route-network validation, retrospective DOI-backed route recovery and route-ranking ablations. Current claims are limited to route-aware screening and retrospective route-family recovery; blind time-split validation, calibrated success probabilities and prospective synthesis tests require populated validation datasets.
% 中文：我们通过能量引擎测试、benchmark 尺度路线网络验证、带 DOI 支撑的回顾性路线恢复以及路线排序消融来评价这一对象。当前主张仅限于路线感知筛选和回顾性路线族恢复；盲法时间切分验证、校准后的成功概率以及前瞻性合成测试都需要填充完成的验证数据集。

\section*{Results}
% 中文：结果

\subsection*{From target labels to route-network decisions}
% 中文：从目标标签到路线网络决策
\begin{figure*}[htbp] 
    \centering
    \includegraphics[width=1.0\linewidth]{figures/Fig1_formula_CIF_to_route_workflow.pdf}
    \caption{
        \textbf{Figure 1 | InvDesRoute defines a route-network decision object.}
        \textbf{(a)} Workflow illustrated with BaTiO$_3$. Formula or CIF input is mapped to the Ba--Ti--O chemical system, expanded with Materials Project phases and practical precursor/reservoir entries, evaluated by PET-OAM-XL and converted into a balanced reaction network. For BaTiO$_3$, 651 formation routes are generated and ranked into synthesis recommendations. The top synthesis-aware routes are BaCO$_3$ + TiO$_2$ $\rightarrow$ BaTiO$_3$ + CO$_2$ (score 0.951) and BaO + TiO$_2$ $\rightarrow$ BaTiO$_3$ (score 0.943), with high target confidence and CO$_2$/finite-temperature refinement flags.
        \textbf{(b)} Route-ranking score components, including precursor/template priors, reaction energy, practical-route status, complexity, phase availability and condition simplicity.
        \textbf{(c)} SrTiO$_3$ example showing target-scale validation: 25 phases, 673 routes, PET/MP-form MAE of 0.089 eV per target, 97\% sign consistency and exact top-1 recovery of the literature carbonate route SrCO$_3$ + TiO$_2$ $\rightarrow$ SrTiO$_3$ + CO$_2$.}
        %The workflow begins from either a chemical formula or a user-supplied CIF for the target phase, illustrated here with BaTiO$3$. The target is first mapped onto a chemical-system-aware search space, after which the Materials Project phase pool is expanded with stable phases and a library of experimentally common precursors. Representative phases and precursor structures are passed through an atomic-environment representation and a machine-learning potential to obtain consistent phase energies. These energies are then used to enumerate balanced formation reactions across the expanded phase network, including elemental, oxide, carbonate and gas-mediated hypotheses, and to compute reaction energies on a common thermodynamic scale. The resulting route network is finally ranked by a synthesizability layer that combines reaction energy, target phase stability ($E\mathrm{hull}$), precursor priors, route-template priors, route complexity, gas involvement, competing decomposition risk and refinement flags. The output is a ranked synthesis report that identifies the most actionable route hypotheses rather than a binary synthesis label.
    }
    \label{fig:workflow}
\end{figure*}
\invdesroute{} defines the screening unit as a route-network record rather than a target label. For a formula or CIF input, the record links the target definition to a Materials Project phase context, common precursor and gas reservoirs, PET-OAM-XL phase energies, balanced formation and decomposition reactions, route ranks and refinement flags. The same object therefore carries both the candidate synthesis hypotheses and the evidence needed to decide which hypotheses should be inspected first.
% 中文：\invdesroute{} 将筛选单元定义为路线网络记录，而不是目标标签。对于化学式或 CIF 输入，该记录把目标定义与 Materials Project 相背景、常见前驱体和气体库、PET-OAM-XL 相能量、配平的形成与分解反应、路线排名和精修标记连接起来。因此，同一个对象同时承载候选合成假设，以及判断哪些假设应优先检查所需的证据。


This design separates two quantities that are often conflated in synthesizability screens. The complete network is retained as a thermodynamic search space, while the displayed recommendation window is a route-prioritized view of that space. Each displayed route is accompanied by interpretable evidence---reaction favorability, target energy above hull, precursor class, route template, complexity, gas involvement, decomposition competition and model-disagreement flags---so a high rank is presented as a checkable synthesis hypothesis, not as a stand-alone success label.
% 中文：这种设计区分了可合成性筛选中经常被混淆的两个量。完整网络作为热力学搜索空间被保留，而展示出来的推荐窗口则是该空间中按路线优先级排序的视图。每条展示路线都带有可解释证据，包括反应有利性、目标 energy above hull、前驱体类别、路线模板、复杂度、气体参与、分解竞争和模型不一致标记，因此高排名被呈现为可检查的合成假设，而不是独立的成功标签。



\noindent Figure~\ref{fig:workflow} defines the route-network object; the next four figures test whether it can be scored, generated at scale, compared with literature chemistry and prioritized for refinement. Workflow modules are summarized in Supplementary Table 1.
% 中文：Figure~\ref{fig:workflow} 定义了路线网络对象；接下来的四个图检验它是否可以被评分、规模化生成、与文献化学比较，并按精修优先级排序。工作流模块汇总于 Supplementary Table 1。

\subsection*{Route-scale energy consistency selects the energy engine}
% 中文：路线尺度能量一致性决定能量引擎

Route networks amplify energy errors differently from single-structure benchmarks: a route energy is a stoichiometric difference across reactants, target phases and byproducts. We therefore evaluated candidate machine-learning interatomic potentials from the current universal-potential landscape \cite{chen2022m3gnet,deng2023chgnet,batatia2022mace,park2024sevennet,yang2024mattersim,pozdnyakov2023pet,mazitov2025petmad,bigi2026upet} on a DFT-referenced complex-hydride benchmark with 107 crystalline compounds, 416 balanced reactions and 37 target hydrides. The compound, reaction and target levels test whether an energy model preserves the scale, sign and ordering needed for route-network screening.
% 中文：路线网络放大能量误差的方式不同于单结构 benchmark：路线能量是反应物、目标相和副产物之间按化学计量相减得到的差值。因此，我们在一个以 DFT 为参考的复杂氢化物 benchmark 上，评价了当前通用势模型体系中的候选机器学习原子间势 \cite{chen2022m3gnet,deng2023chgnet,batatia2022mace,park2024sevennet,yang2024mattersim,pozdnyakov2023pet,mazitov2025petmad,bigi2026upet}；该 benchmark 包含 107 个晶体化合物、416 条配平反应和 37 个目标氢化物。化合物、反应和目标三个层级共同检验能量模型是否能保持路线网络筛选所需的尺度、符号和排序。

PET-OAM-XL gives the strongest route-relevant performance across these levels. It reaches a phase-energy MAE of 0.0061 eV atom$^{-1}$, a reaction-energy MAE of 0.110 eV and RMSE of 0.165 eV across 416 reactions, 98.8\% reaction-sign consistency, and reproduces the DFT-ranked top route for all 37 hydride targets. A Li$_2$AgH$_6$/Li$_2$AuH$_6$ pressure stress test, motivated by recent hydride-computing studies and the sensitivity of these candidates to structural and dynamical stability checks \cite{ouyang2025li2auh6,gao2025maximum_tc,wei2026li2auh6_pressure,ding2026kinetic_instability}, gives the same qualitative outcome: across 68 reactions at 0 and 10 GPa, PET-OAM-XL reaches a MAE of 0.249 eV and preserves all reaction signs. These results support using PET-OAM-XL as the production energy engine for the current route-screening tests before evaluating the complete workflow beyond hydride chemistry.
% 中文：PET-OAM-XL 在这些层级上给出了与路线最相关的最强表现。它达到 0.0061 eV atom$^{-1}$ 的相能量 MAE，在 416 条反应上达到 0.110 eV 的反应能 MAE 和 0.165 eV 的 RMSE，反应符号一致率为 98.8\%，并为全部 37 个氢化物目标复现了 DFT 排名第一的路线。受近期氢化物计算研究以及这些候选物对结构和动力学稳定性检查敏感性的启发 \cite{ouyang2025li2auh6,gao2025maximum_tc,wei2026li2auh6_pressure,ding2026kinetic_instability}，Li$_2$AgH$_6$/Li$_2$AuH$_6$ 压力应力测试给出了相同的定性结果：在 0 和 10 GPa 下的 68 条反应中，PET-OAM-XL 达到 0.249 eV 的 MAE，并保持所有反应符号不变。这些结果支持在将完整工作流扩展到氢化物化学之外之前，把 PET-OAM-XL 用作当前路线筛选测试的生产能量引擎。

\begin{figure*}[htbp] 
    \centering
    \includegraphics[width=1.0\linewidth]{figures/Fig2.pdf}
    \caption
    {\textbf{Benchmark evidence and route-level validation of InvDesRoute.}
    \textbf{(a)} PET-OAM-XL shows the best phase-energy, reaction-energy and reaction-sign performance among the tested interatomic potentials, supporting its use as the production energy engine.
    \textbf{(b)} In the 120-target benchmark, InvDesRoute produces 17,085 PET/MP-form comparable routes across five material classes. PET and MP-form references give 93.1\% mean sign consistency, 99/120 top-1 route agreement and 3.68/5 mean top-5 overlap; the recommendation layer further identifies 118/120 exergonic practical routes, 117/120 exergonic low-Ehull cases and 104/15/1 high/medium/low-confidence targets.
    \textbf{(c)} Literature-route validation shows that route-family matching is essential for experimental relevance. The full networks contain moderate-or-better analogues for 143/145 DOI-backed routes, and top-5 recovery improves from 13/145 with PET-energy-only ranking to 113/145 after InvDesRoute scoring and diversity-aware display.}
    \label{fig2}
\end{figure*}

\subsection*{Benchmark-scale route-network validation across 120 targets}
% 中文：跨 120 个目标的 benchmark 尺度路线网络验证

We next tested whether the route-network record can be generated and compared at benchmark scale. The 120-target set spans halides, binary oxides, nitrides/carbides/chalcogenides, ternary oxides and battery, salt or silicate compounds. \invdesroute{} completes all targets, retrieves 1,619 phases, evaluates PET energies for 1,596 CIF-bearing phases, enumerates 18,063 formation reactions and yields 17,085 PET/MP-form comparable reactions. PET coverage is 98.6\%, indicating that most missing comparisons arise from phase availability and filtering rather than route-generation failure.
% 中文：接下来，我们测试路线网络记录是否可以在 benchmark 尺度上生成和比较。120 个目标集合覆盖卤化物、二元氧化物、氮化物/碳化物/硫族化合物、三元氧化物，以及电池、盐或硅酸盐化合物。\invdesroute{} 完成了全部目标，检索 1,619 个相，为 1,596 个含 CIF 的相评估 PET 能量，枚举 18,063 条形成反应，并得到 17,085 条可进行 PET/MP-form 比较的反应。PET 覆盖率为 98.6\%，说明大多数缺失比较来自相可用性和过滤，而不是路线生成失败。

Across targets, PET and MP-form reaction-energy references agree on reaction sign with mean consistency of 0.931. The top-ranked route agrees for 99 of 120 targets, and mean top-5 overlap is 3.68. Halides show full top-1 agreement; ternary oxides and battery/salt/silicate compounds produce denser networks and lower top-1 agreement, consistent with near-degenerate alternatives. The practical-route layer retains chemically legible hypotheses for most targets: 118 of 120 have at least one exergonic practical PET route, while the complete thermodynamic network remains available for inspection.
% 中文：在所有目标上，PET 与 MP-form 反应能参考在反应符号上的平均一致性为 0.931。120 个目标中有 99 个的最高排序路线一致，平均 top-5 重叠数为 3.68。卤化物表现出完全的 top-1 一致性；三元氧化物以及电池/盐/硅酸盐化合物产生更密集的网络和更低的 top-1 一致性，这与近简并替代路线相符。practical-route 层为大多数目标保留了化学上可读的假设：120 个目标中有 118 个至少有一条放热的 practical PET 路线，同时完整热力学网络仍可供检查。

To test portability to experimentally occurring compositions, we added 53 literature-seeded targets and 86 DOI-backed route records absent from the 120-target benchmark. This extension produces 2,769 PET/MP-form comparable reactions, retains a mean sign consistency of 0.904 and agrees with the MP-form top-ranked route for 42 of 53 targets. Combining both sets gives a 173-target thermodynamic validation set with 19,854 comparable reactions, mean sign consistency of 0.923 and top-1 agreement for 141 targets.
% 中文：为测试对实验中出现的组成的可迁移性，我们加入了 53 个文献种子目标和 86 条不在 120 目标 benchmark 中的 DOI 支撑路线记录。该扩展产生 2,769 条可进行 PET/MP-form 比较的反应，平均符号一致性保持在 0.904，并在 53 个目标中的 42 个上与 MP-form 最高排序路线一致。合并两个集合后得到一个包含 173 个目标的热力学验证集，其中有 19,854 条可比较反应，平均符号一致性为 0.923，141 个目标具有 top-1 一致性。



\noindent Category statistics, literature-seeded statistics, representative routes and claim boundaries are summarized in Supplementary Tables 2--4 and 14.
% 中文：类别统计、文献种子统计、代表性路线和主张边界汇总于 Supplementary Tables 2--4 和 14。

\subsection*{Literature-route recovery tests route chemistry}
% 中文：文献路线恢复检验路线化学

Internal thermodynamic consistency is necessary but not sufficient: a route-network screen should also recover the chemistry that experimental papers report. We therefore compared generated networks with DOI-backed solid-state synthesis records from the Kononova \textit{et al.} text-mined corpus \cite{kim2017synthesis_insights,kononova2019textmined,wang2022solution_synthesis}. Exact reaction identity is intentionally strict but brittle; route-family similarity instead records component-wise agreement in precursor classes, target-element sources, volatile products and synthesis templates. This uses literature routes as retrospective chemistry anchors while recognizing reporting and historical biases in extracted corpora \cite{jia2019anthropogenic_biases}. The component score builds on evidence that precursor similarity carries synthesis information \cite{he2020precursor_similarity,szymanski2023arrows3} and allows carbonate, hydroxide/hydrate, ammonium-phosphate and gas-mediated routes to be compared with computed route families.
% 中文：内部热力学一致性是必要的，但并不充分：路线网络筛选还应恢复实验论文所报道的化学路线。因此，我们将生成的网络与 Kononova \textit{et al.} 文本挖掘语料库中的带 DOI 支撑固相合成记录进行比较 \cite{kim2017synthesis_insights,kononova2019textmined,wang2022solution_synthesis}。精确反应恒等性被有意设定得严格，但也很脆弱；路线族相似性则记录前驱体类别、目标元素来源、挥发性产物和合成模板中的逐组分一致性。这把文献路线作为回顾性化学锚点，同时承认提取语料中的报道偏差和历史偏差 \cite{jia2019anthropogenic_biases}。组分评分建立在前驱体相似性携带合成信息的证据之上 \cite{he2020precursor_similarity,szymanski2023arrows3}，并允许将碳酸盐、氢氧化物/水合物、铵盐磷酸盐和气体介导路线与计算路线族进行比较。

A curated 80-route seed set verifies route availability: all 80 seed reactions appear in the full network, 75 appear in the top-5 under strict identity and all 80 appear in the top-20 under route-family matching. Across 145 DOI-backed routes, the full networks contain strong-or-exact route-family analogues for 141 routes and moderate-or-better analogues for 143. We report these route-family results as a three-level evidence ladder. The current evidence is an algorithmic label computed from frozen component scores, rule versions and thresholds and remains checkable through the written literature precursors, matched computed route, similarity components and route ranks. A stronger robustness claim requires threshold perturbations and negative-control route pairs. A human-supported claim additionally requires blinded annotation packets that hide model score, model rank, expected labels and algorithmic labels until independent labels and agreement statistics are recorded.
% 中文：一个经过整理的 80 条路线种子集用于验证路线可用性：全部 80 条种子反应都出现在完整网络中，75 条在严格恒等标准下出现在 top-5 中，全部 80 条在路线族匹配下出现在 top-20 中。在 145 条带 DOI 支撑的路线中，完整网络为 141 条路线包含 strong-or-exact 的路线族类似物，并为 143 条路线包含 moderate-or-better 的类似物。我们将这些路线族结果报告为三层证据阶梯。当前证据是由冻结的组分评分、规则版本和阈值计算得到的算法标签，并可通过文献中写明的前驱体、匹配的计算路线、相似性组分和路线排名进行检查。更强的稳健性主张需要阈值扰动和负对照路线对。人工支持的主张还需要盲法标注包，在记录独立标签和一致性统计之前隐藏模型分数、模型排名、预期标签和算法标签。



\noindent Literature-route validation is summarized in Supplementary Table 5.
% 中文：文献路线验证汇总于 Supplementary Table 5。

\subsection*{\invdesroute{} ranking makes useful routes visible}
% 中文：\invdesroute{} 排序使有用路线可见

The literature comparison separates route existence from route visibility. To isolate the ranking layer, we compared four display policies on the same 145 DOI-backed route records while holding PET energies, generated candidate routes and route-family labels fixed: PET-energy-only sorting, practical-first energy sorting, the full \invdesroute{} score and diversity-aware display. This is a matched route-level ablation, not a claim that route-aware ranking outperforms target-level stability screens. That workflow-level claim is reserved for frozen baselines with shared DOI routes, candidate routes and top-$k$ windows, pre-label ranks, an explicit target-to-route exposure source and no literature-label inputs.
% 中文：文献比较把路线存在性与路线可见性区分开来。为隔离排序层，我们在保持 PET 能量、生成的候选路线和路线族标签固定的条件下，对同一批 145 条带 DOI 支撑的路线记录比较了四种展示策略：仅按 PET 能量排序、practical-first 能量排序、完整 \invdesroute{} 评分以及 diversity-aware 展示。这是匹配的路线级消融，并不声称路线感知排序优于目标级稳定性筛选。那种工作流级主张需要冻结的基线，并共享 DOI 路线、候选路线和 top-$k$ 窗口，包含标签前排名、明确的目标到路线暴露来源，且评分输入中不包含文献标签。

Within this fixed route set, energy-only ranking is a weak proxy for experimental route relevance. In the top-5 window, pure PET-energy sorting recovers only 13 of 145 moderate-or-better literature-route families, and practical-first sorting recovers 33. The full \invdesroute{} score recovers 109 moderate-or-better and 87 strong-or-exact families at top-5; diversity-aware display increases moderate-or-better top-5 recovery to 113 and top-10 recovery to 126. The recommendation layer changes which hypotheses become visible by combining reaction favorability with precursor priors, route templates, complexity, phase availability and condition-risk terms.
% 中文：在这个固定路线集合中，仅按能量排序是实验路线相关性的弱代理。在 top-5 窗口中，纯 PET 能量排序只能恢复 145 个 moderate-or-better 文献路线族中的 13 个，practical-first 排序恢复 33 个。完整 \invdesroute{} 评分在 top-5 中恢复 109 个 moderate-or-better 和 87 个 strong-or-exact 路线族；diversity-aware 展示进一步将 moderate-or-better 的 top-5 恢复数提高到 113，将 top-10 恢复数提高到 126。推荐层通过结合反应有利性、前驱体先验、路线模板、复杂度、相可用性和条件风险项，改变了哪些假设能够变得可见。

The remaining failures are chemically organized. Binary-oxide and carbonate-source routes are usually present and highly ranked, whereas hydroxide/hydrate, nitrate/ammonium, phosphate and oxalate/organic-salt routes are often present in the full network but absent from the top-5 window. These classes use written precursor identity as a proxy for dehydration, decomposition, volatilization, redox chemistry or gas release, pointing to precursor-decomposition thermochemistry and finite-temperature reservoir terms as targeted upgrades.
% 中文：剩余失败具有清晰的化学组织性。二元氧化物和碳酸盐来源路线通常存在且排名较高，而氢氧化物/水合物、硝酸盐/铵盐、磷酸盐以及草酸盐/有机盐路线往往存在于完整网络中，却没有进入 top-5 窗口。这些类别使用文献写明的前驱体身份作为脱水、分解、挥发、氧化还原化学或放气过程的代理，指向前驱体分解热化学和有限温库项这两个有针对性的升级方向。



\noindent Ranking, baseline-design, precursor-diagnostic and recommendation summaries are provided in Supplementary Tables 6--9.
% 中文：排序、基线设计、前驱体诊断和推荐汇总见 Supplementary Tables 6--9。

\subsection*{Uncertainty labels define a refinement queue}
% 中文：不确定性标签定义精修队列

Confidence labels in \invdesroute{} are deterministic, trigger-based refinement labels, not success probabilities. They combine observable route and phase evidence--exergonic practical-route support, target energy above hull, PET/MP-form top-route agreement, top-5 overlap and decomposition competition--to decide what should be refined next. In the 120-target benchmark, 26 targets enter the DFT/expert-refinement queue, led by PET/MP-form top-1 route mismatch in 21 targets. Nine targets show exergonic decomposition competition, two lack an exergonic practical PET route and one receives a low-confidence label.
% 中文：\invdesroute{} 中的置信标签是确定性的、基于触发条件的精修标签，而不是成功概率。它们结合可观测的路线和相证据，包括放热 practical-route 支撑、目标 energy above hull、PET/MP-form 最高路线一致性、top-5 重叠以及分解竞争，以决定下一步应精修什么。在 120 目标 benchmark 中，26 个目标进入 DFT/专家精修队列，其中主要原因是 21 个目标存在 PET/MP-form top-1 路线不匹配。9 个目标显示放热分解竞争，2 个缺少放热 practical PET 路线，1 个获得低置信标签。

Examples clarify the labels. CoO is the only low-confidence target: its best practical formation route is nearly thermoneutral but slightly endergonic by PET, PET and MP-form select different top-ranked routes and decomposition competition is flagged. InN instead has PET/MP-form top-route agreement but a slightly endergonic best practical elemental route and decomposition competition. Thus the label names should be read as action priorities--which calculation, thermochemical datum or expert review is most likely to change the recommendation--not as calibrated likelihoods of experimental success.
% 中文：具体例子可澄清这些标签的含义。CoO 是唯一低置信目标：其最佳 practical 形成路线接近热中性，但按 PET 略微吸热；PET 和 MP-form 选择不同的最高排序路线，并且存在分解竞争标记。相比之下，InN 具有 PET/MP-form 最高路线一致性，但其最佳 practical 元素路线略微吸热，同时也存在分解竞争。因此，这些标签名称应被理解为行动优先级，即哪项计算、热化学数据或专家审查最可能改变推荐，而不是实验成功可能性的校准概率。

\section*{Discussion}
% 中文：讨论

The central result is that \invdesroute{} turns inorganic synthesis screening from target-stability classification into route-network decision making. Its contribution is not another stability score: it builds an evidence-bearing route object in which structure energetics, balanced reaction networks, phase competition, precursor chemistry, ranking choices and refinement needs can be inspected together. The benchmark and literature analyses show that this object can be generated at scale, compared against thermodynamic references and used to make experimentally familiar route families visible without hiding the lower-level route network.
% 中文：核心结果是，\invdesroute{} 将无机合成筛选从目标稳定性分类转变为路线网络决策。它的贡献并不是另一个稳定性分数，而是构建了一个承载证据的路线对象，使结构能量学、配平反应网络、相竞争、前驱体化学、排序选择和精修需求能够被共同检查。benchmark 和文献分析表明，这一对象可以规模化生成，可以与热力学参考进行比较，并能在不隐藏底层路线网络的前提下，使实验上熟悉的路线族变得可见。

The limitations define the next validation experiments. The literature comparison is retrospective and inherits the coverage and reporting biases of published synthesis records \cite{jia2019anthropogenic_biases}; blind-style recovery requires ranks to be submitted from a frozen scoring, candidate-generation and precursor-library state before held-out DOI labels are inspected \cite{kapoor2023leakage}. Route-family similarity should be stress-tested with threshold perturbations, negative controls and blinded annotation, while confidence labels remain triage categories until held-out outcomes support calibration diagnostics \cite{guo2017calibration,gneiting2007proper}. Hydroxide/hydrate, nitrate/ammonium, phosphate and oxalate failures point to missing precursor-decomposition and volatile-reservoir thermochemistry rather than missing enumeration.
% 中文：这些局限性定义了下一步验证实验。文献比较是回顾性的，并继承了已发表合成记录的覆盖偏差和报道偏差 \cite{jia2019anthropogenic_biases}；类盲法恢复要求在检查 held-out DOI 标签之前，从冻结的评分、候选生成和前驱体库状态提交排名 \cite{kapoor2023leakage}。路线族相似性应通过阈值扰动、负对照和盲法标注进行压力测试，而置信标签在 held-out 结果支持校准诊断之前仍应保留为 triage 类别 \cite{guo2017calibration,gneiting2007proper}。氢氧化物/水合物、硝酸盐/铵盐、磷酸盐和草酸盐相关失败指向缺失的前驱体分解与挥发性库热化学，而不是枚举缺失。

\invdesroute{} therefore reframes inorganic synthesizability as an inspectable route-network decision: a promising target is one whose network contains favorable, chemically interpretable formation routes, tolerable phase competition and explicit uncertainties for refinement before experimental or higher-fidelity computational effort is spent.
% 中文：因此，\invdesroute{} 将无机可合成性重新表述为一种可检查的路线网络决策：一个有前景的目标，是其网络在投入实验或更高保真计算之前，已经包含有利且化学上可解释的形成路线、可容忍的相竞争，以及明确需要精修的不确定性。

\section*{Methods}
% 中文：方法

\subsection*{Phase retrieval and target handling}
% 中文：相检索和目标处理

\invdesroute{} accepts a formula or CIF target. Formula inputs are reduced with \texttt{pymatgen} \cite{ong2013pymatgen}; for target elements $\mathcal{E}$, all non-empty sub-chemical systems are queried in Materials Project to retrieve identifiers, reduced formulas, structures, energies, formation energies and energies above hull \cite{jain2013materialsproject}. Retrieved structures and phase metadata are cached for energy evaluation and reaction enumeration.
% 中文：\invdesroute{} 接受化学式或 CIF 目标。化学式输入使用 \texttt{pymatgen} \cite{ong2013pymatgen} 约化；对于目标元素集合 $\mathcal{E}$，工作流会在 Materials Project 中查询所有非空子化学体系，以检索标识符、约化化学式、结构、能量、形成能和 energy above hull \cite{jain2013materialsproject}。检索得到的结构和相元数据会被缓存，用于后续能量评估和反应枚举。

For benchmark calculations, entries were sorted by energy above hull and deduplicated by normalized reduced formula, retaining the lowest-Ehull representative while keeping full phase metadata for later polymorph-sensitive refinement. For CIF input, the uploaded structure is parsed directly, its composition is inferred from the atomic contents and the user-provided structure is inserted as the target phase before enumeration. This enables screening of metastable, hypothetical or experimentally proposed structures that are not Materials Project ground-state representatives.
% 中文：在 benchmark 计算中，条目按 energy above hull 排序，并根据规范化后的约化化学式去重，保留最低 Ehull 代表，同时保留完整相元数据以供后续对多晶型敏感的精修。对于 CIF 输入，上传结构会被直接解析，其组成由原子内容推断，并在枚举前将用户提供的结构作为目标相插入。这使工作流能够筛选不属于 Materials Project 基态代表的亚稳、假想或实验提出结构。

For literature-route validation and recommendations, the phase pool is augmented with common elemental reservoirs and precursor classes, including oxides, hydroxides, carbonates, nitrates, oxalates, ammonium phosphates, hydrated salts and simple gas reservoirs. Species with reliable Materials Project representatives use the same MP metadata and PET energies as other phases. Volatile reservoirs or process-level precursors without direct MP solid representatives are retained as labelled route objects with availability flags and, when applicable, approximate reservoir references or gas-chemical-potential corrections.
% 中文：对于文献路线验证和推荐，相池会加入常见元素库和前驱体类别，包括氧化物、氢氧化物、碳酸盐、硝酸盐、草酸盐、铵盐磷酸盐、水合盐和简单气体库。具有可靠 Materials Project 代表的物种使用与其他相相同的 MP 元数据和 PET 能量。没有直接 MP 固相代表的挥发性库物种或过程级前驱体，则作为带标签的路线对象保留，并带有可用性标记；适用时还包括近似库参考或气体化学势校正。

\subsection*{PET-OAM-XL energy evaluation}
% 中文：PET-OAM-XL 能量评估

PET-OAM-XL was used as the production energy model for route ranking. Cached CIF structures were converted with ASE and evaluated in the uPET/PET environment \cite{hjorthlarsen2017ase,pozdnyakov2023pet,mazitov2025petmad,bigi2026upet}. Each unique structure was evaluated once for a fixed model configuration, and the predicted energy was stored with phase identity, formula, structural source and thermodynamic metadata.
% 中文：PET-OAM-XL 被用作路线排序的生产能量模型。缓存的 CIF 结构通过 ASE 转换，并在 uPET/PET 环境中评估 \cite{hjorthlarsen2017ase,pozdnyakov2023pet,mazitov2025petmad,bigi2026upet}。每个唯一结构在固定模型配置下只评估一次，预测能量会与相身份、化学式、结构来源和热力学元数据一起存储。

Reaction screening depends on cross-phase energy consistency rather than isolated phase accuracy alone. We therefore evaluate PET-OAM-XL at both phase-energy and reaction-energy levels; in the workflow output, PET-OAM-XL ranks routes and Materials Project formation energies provide an internal consistency reference.
% 中文：反应筛选依赖跨相能量一致性，而不仅是孤立相精度。因此，我们同时在相能量和反应能层面评价 PET-OAM-XL；在工作流输出中，PET-OAM-XL 用于路线排序，而 Materials Project 形成能提供内部一致性参考。

\subsection*{Reaction enumeration and energy normalization}
% 中文：反应枚举和能量归一化

Formation reactions place the target on the product side and enumerate non-target reactant sets, with optional byproducts also allowed. Decomposition reactions place the target on the reactant side and enumerate product sets from the same pool. Integer element balance is solved for each candidate; reactions with invalid target coefficients, unstable balances or coefficient scales above the enumeration limits are removed. Each retained reaction stores the balanced expression, formulas, target coefficient, PET and \mpform{} reaction energies, practical-route descriptors and complexity features.
% 中文：形成反应将目标放在产物侧，并枚举非目标反应物集合，同时也允许可选副产物。分解反应将目标放在反应物侧，并从同一相池中枚举产物集合。每个候选反应都会求解整数元素配平；目标系数无效、配平不稳定或系数尺度超过枚举限制的反应会被移除。每条保留反应都会存储配平表达式、化学式、目标系数、PET 和 \mpform{} 反应能、practical-route 描述符以及复杂度特征。

The Materials Project reference used in the final workflow is based on formation energies. For each phase, the MP formation energy per atom is converted to a per-formula-unit value by multiplying by the number of atoms in the reduced formula,
% 中文：最终工作流使用的 Materials Project 参考基于形成能。对于每个相，MP 每原子形成能通过乘以约化化学式中的原子数，转换为每化学式单位的值，
\begin{equation}
E^{\mathrm{form}}_{\mathrm{fu}} = E^{\mathrm{form}}_{\mathrm{atom}} N_{\mathrm{atom,fu}}.
\end{equation}

The \mpform{} reaction energy is then evaluated on the same balanced reaction expression used for PET,
% 中文：随后，\mpform{} 反应能在与 PET 相同的配平反应表达式上计算，
\begin{equation}
\Delta E^{\mathrm{MP-form}}_{\mathrm{rxn}}
=
\sum_{i \in \mathrm{products}} \nu_i E^{\mathrm{form}}_{i,\mathrm{fu}}
-
\sum_{j \in \mathrm{reactants}} \nu_j E^{\mathrm{form}}_{j,\mathrm{fu}},
\end{equation}
where $\nu_i$ and $\nu_j$ are positive stoichiometric coefficients. PET reaction energies use the same convention, and all reaction energies are normalized by the target coefficient unless otherwise specified. Direct Materials Project total-energy subtraction is not used as the primary reference because it differs from formation-energy differences and is not directly comparable to energy-above-hull descriptors. The route graph is a screening object: it records balanced formation and decomposition hypotheses and flags cases needing condition-dependent refinement, while leaving phonon free energies, chemical-potential windows, kinetic barriers and phase-evolution simulations to downstream checks \cite{buckeridge2014chemical_potentials,togo2015phonon,henkelman2000neb,gallant2024reactca}.
% 中文：其中 $\nu_i$ 和 $\nu_j$ 是正化学计量系数。PET 反应能使用相同约定；除非另有说明，所有反应能都按目标系数归一化。直接相减 Materials Project 总能量不作为主要参考，因为它不同于形成能差，也不能直接与 energy-above-hull 描述符比较。路线图是一个筛选对象：它记录配平的形成和分解假设，并标记需要条件依赖精修的情况，而将声子自由能、化学势窗口、动力学势垒和相演化模拟留给下游检查 \cite{buckeridge2014chemical_potentials,togo2015phonon,henkelman2000neb,gallant2024reactca}。

\subsection*{Practical-route filters and ranking metrics}
% 中文：Practical-route 过滤器和排序指标

The complete network was retained for thermodynamic analysis; chemically legible first recommendations used a practical-route subset. A route enters this subset when reduced non-target formulas contain at most 10 atoms, the maximum-to-minimum nonzero coefficient ratio is no larger than 5, non-target Ehull is no larger than 0.10 eV atom$^{-1}$ when available, and the total coefficient scale per target is no larger than 30. These thresholds define a reproducible display filter and do not remove reactions from the full network.
% 中文：完整网络被保留用于热力学分析；化学上可读的首批推荐使用 practical-route 子集。当约化后的非目标化学式最多含 10 个原子、非零系数的最大/最小比不大于 5、非目标 Ehull 在可用时不大于 0.10 eV atom$^{-1}$，并且每个目标的总系数尺度不大于 30 时，路线进入该子集。这些阈值定义了可复现的展示过滤器，但不会从完整网络中移除反应。

PET-OAM-XL and \mpform{} were compared on identical enumerated formation-reaction sets using target-level MAE, RMSE, signed bias, sign consistency, top-1 route agreement and top-5 route overlap. Sign consistency tests whether the reaction sign is preserved; top-1 agreement requires identical lowest-energy reaction expressions; top-5 overlap counts shared reaction expressions among the five lowest-energy PET and \mpform{} routes. These metrics test route ordering, not only average reaction-energy magnitude.
% 中文：PET-OAM-XL 与 \mpform{} 在相同的枚举形成反应集合上比较，指标包括目标级 MAE、RMSE、有符号偏差、符号一致性、top-1 路线一致性和 top-5 路线重叠。符号一致性检验反应符号是否被保持；top-1 一致性要求最低能反应表达式完全相同；top-5 重叠统计 PET 和 \mpform{} 五条最低能路线中共享的反应表达式数量。这些指标检验路线排序，而不仅是平均反应能大小。

\subsection*{Recommendation labels and refinement queue}
% 中文：推荐标签和精修队列

Target-level confidence labels are deterministic refinement-triage categories. They combine exergonic practical-route support, target energy above hull, PET/\mpform{} top-route agreement, PET/\mpform{} top-5 overlap and competing decomposition strength. High-confidence targets have favorable practical PET routes, low target Ehull and no major disagreement or decomposition flags; medium-confidence targets retain plausible support but require refinement of at least one uncertainty source; low-confidence targets lack support in the screened static network. The label generator stores the trigger fields and row-level label identifiers so that any later calibration table can join back to the unchanged pre-outcome label.
% 中文：目标级置信标签是确定性的精修分流类别。它们结合放热 practical-route 支撑、目标 energy above hull、PET/\mpform{} 最高路线一致性、PET/\mpform{} top-5 重叠以及竞争分解强度。高置信目标具有有利的 practical PET 路线、较低目标 Ehull，且没有主要不一致或分解标记；中置信目标仍有合理支撑，但至少一个不确定性来源需要精修；低置信目标在筛选得到的静态网络中缺少支撑。标签生成器保存触发字段和行级标签标识符，使后续校准表可以回连到未改变的 pre-outcome 标签。

The same flags define a DFT/expert-refinement queue, linking model-sensitive rankings to DFT checks, gas-sensitive routes to thermochemical correction, decomposition-prone targets to phase-diagram refinement and precursor-sensitive routes to expert review or targeted precursor thermochemistry.
% 中文：同一组标记定义 DFT/专家精修队列，将模型敏感排序连接到 DFT 检查，将气体敏感路线连接到热化学校正，将易分解目标连接到相图精修，并将前驱体敏感路线连接到专家审查或定向前驱体热化学。

Calibration is a separate held-out outcome-validation task, not part of the current 120-target benchmark. A probability claim requires frozen calibration rules, trigger fields, outcome definitions, split policy, freeze manifest, score-bin policy and predeclared label-wise support thresholds before held-out outcomes are inspected. Populated calibration rows must join to the frozen confidence-label record without changing the label, row identity or trigger-field values. Only then can label-wise reliability, expected calibration error or predicted success probabilities be reported. Until those records are populated and the minimum records per label are met, high, medium and low remain refinement-triage labels rather than calibrated probabilities.
% 中文：校准是独立的 held-out 结果验证任务，不属于当前 120 目标 benchmark。任何概率主张都要求在查看 held-out 结果之前冻结校准规则、触发字段、结果定义、划分策略、freeze manifest、分数分箱策略和预先声明的逐标签支持阈值。填充后的校准行必须回连到冻结的置信标签记录，且不能改变标签、行身份或触发字段值。只有这样才可以报告逐标签可靠性、expected calibration error 或预测成功概率。在这些记录被填充且每类标签达到最低记录数之前，high、medium 和 low 仍是精修分流标签，而不是校准概率。

\subsection*{Literature-route extraction and route-family similarity}
% 中文：文献路线提取和路线族相似性

Literature validation used the public text-mined solid-state synthesis dataset of Kononova \textit{et al.} and its open repository \cite{kononova2019textmined,cedergrouphub_textmined_synthesis_public}. Parsed recipes were filtered to retain solid-state entries with valid target formulas, identifiable precursor formulas, DOI provenance and target compositions on the product side; doped compositions, composites, variable formulas, ambiguous mixtures and targets outside the validation set were excluded. This yielded 59 DOI-backed routes from 31,251 parsed solid-state recipes for the initial 80-target comparison. The literature routes were used only for external validation, not to train PET-OAM-XL or fit reaction energies.
% 中文：文献验证使用 Kononova \textit{et al.} 的公开文本挖掘固相合成数据集及其开放仓库 \cite{kononova2019textmined,cedergrouphub_textmined_synthesis_public}。解析后的配方经过过滤，只保留具有有效目标化学式、可识别前驱体化学式、DOI 来源且目标组成位于产物侧的固相条目；掺杂组成、复合物、可变化学式、模糊混合物以及验证集之外的目标被排除。该流程从 31,251 条解析固相配方中为初始 80 目标比较得到 59 条带 DOI 支撑的路线。文献路线仅用于外部验证，而不用于训练 PET-OAM-XL 或拟合反应能。

Computed and literature routes were compared by exact identity and route-family similarity. Exact matching requires identical normalized precursor and product formula sets after balancing. Route-family similarity captures synthesis-motif agreement rather than identical conditions, including hydrated/dehydrated salts, nitrate or carbonate decomposition, ammonium phosphate chemistry and gas-releasing transformations. The score combines precursor overlap, target-element-source recall, precursor-class overlap, volatile/gas terms, element overlap and template bonuses:
% 中文：计算路线和文献路线通过精确恒等性和路线族相似性进行比较。精确匹配要求配平后具有相同的规范化前驱体和产物化学式集合。路线族相似性捕捉的是合成 motif 的一致性，而不是相同条件，包括水合/脱水盐、硝酸盐或碳酸盐分解、铵盐磷酸盐化学以及放气转化。该评分结合前驱体重叠、目标元素来源召回、前驱体类别重叠、挥发物/气体项、元素重叠和模板加分：
\begin{equation}
S_{\mathrm{fam}}
=0.28R_{\mathrm{exact}}
+0.18J_{\mathrm{exact}}
+0.20R_{\mathrm{elem}}
+0.16S_{\mathrm{class}}
+0.08R_{\mathrm{byprod}}
+0.05S_{\mathrm{gas}}
+0.05J_{\mathrm{elem}}
+B_{\mathrm{template}} .
\end{equation}

The terms denote exact precursor recall ($R_{\mathrm{exact}}$), exact precursor Jaccard overlap ($J_{\mathrm{exact}}$), target-element-source recall ($R_{\mathrm{elem}}$), precursor-class overlap ($S_{\mathrm{class}}$), volatile-byproduct recall ($R_{\mathrm{byprod}}$), gas-class agreement ($S_{\mathrm{gas}}$), element overlap ($J_{\mathrm{elem}}$) and template bonuses.
% 中文：这些项分别表示精确前驱体召回 ($R_{\mathrm{exact}}$)、精确前驱体 Jaccard 重叠 ($J_{\mathrm{exact}}$)、目标元素来源召回 ($R_{\mathrm{elem}}$)、前驱体类别重叠 ($S_{\mathrm{class}}$)、挥发性副产物召回 ($R_{\mathrm{byprod}}$)、气体类别一致性 ($S_{\mathrm{gas}}$)、元素重叠 ($J_{\mathrm{elem}}$) 和模板加分。

Template bonuses cover carbonate-source, ammonium-phosphate, hydroxide/hydrate-derived and oxalate-assisted correspondences. Labels are exact, strong, moderate, same-element or weak; exact/strong denote close algorithmic correspondence, whereas moderate denotes recovery of the main target-element sources or precursor classes. Moderate-or-better counts are retrospective recovery under frozen rules, not threshold-independent chemical truth. The release package separates three claim levels: recomputed algorithmic labels; robustness labels after threshold perturbation and negative controls; and human-supported labels after blinded annotation packets, agreement statistics and adjudication records.
% 中文：模板加分覆盖碳酸盐来源、铵盐磷酸盐、氢氧化物/水合物衍生和草酸盐辅助对应关系。标签分为 exact、strong、moderate、same-element 或 weak；exact/strong 表示紧密的算法对应，而 moderate 表示恢复了主要目标元素来源或前驱体类别。moderate-or-better 计数是在冻结规则下的回顾性恢复，而不是不依赖阈值的化学真值。发布包区分三层主张：重新计算得到的算法标签；经过阈值扰动和负对照后的稳健性标签；以及经盲法标注包、一致性统计和裁决记录支持的人工标签。

For the human-supported tier, the blinded annotation layer freezes the annotation unit, visible chemistry fields, hidden score/rank/expected-label/algorithmic-label fields, sampling strata and agreement metrics before interpretation. Agreement records must join to DOI routes and similarity tables and report pairwise agreement, Cohen's kappa \cite{cohen1960agreement}, majority and unanimous labels, adjudicated labels and human-versus-algorithm summaries.
% 中文：在人工支持层，盲法标注层会在解释前冻结标注单元、可见化学字段、隐藏的分数/排名/预期标签/算法标签字段、抽样分层和一致性指标。一致性记录必须连接到 DOI 路线和相似性表，并报告成对一致性、Cohen's kappa \cite{cohen1960agreement}、多数标签和全体一致标签、裁决标签以及人工与算法对比汇总。

Failure-mode diagnostics used the combined 145-route DOI-backed comparison. Literature precursors were assigned to overlapping classes using the route-family rules and summarized by full-network moderate-or-better recovery, top-5 recovery, top-20 recovery and mean top-5 similarity. A diagnostic queue records cases where a moderate-or-better analogue exists in the full network but not in the top-5 window.
% 中文：失败模式诊断使用合并后的 145 条带 DOI 支撑路线比较。文献前驱体使用路线族规则分配到可重叠类别，并按完整网络 moderate-or-better 恢复、top-5 恢复、top-20 恢复和平均 top-5 相似性进行汇总。诊断队列记录完整网络中存在 moderate-or-better 类似物但未进入 top-5 窗口的案例。

\subsection*{Blind-style validation upgrade and leakage controls}
% 中文：类盲法验证升级和泄漏控制

The present literature-route analysis is retrospective and is not used to claim blind generalization. A blind-style or time-split upgrade has three required guards \cite{kapoor2023leakage}. First, scoring, candidate generation, precursor libraries, route-family thresholds, split rules and submitted ranking tables must be frozen before held-out DOI labels are inspected. Second, each held-out row must join to the archived submitted rank in \texttt{frozen\_route\_rankings.csv}, with \texttt{matched\_route\_id} mapped to the frozen ranking \texttt{route\_id} and rank-submission time preceding label release. Third, literature-route, route-family, best-top-$k$, expected, human and adjudicated labels must be excluded from scoring inputs. Until populated frozen-ranking and held-out route tables pass these checks, blind-style recovery remains a validation upgrade rather than numerical evidence.
% 中文：当前文献路线分析是回顾性的，并不用于声称盲法泛化。类盲法或时间切分升级需要三个保护条件 \cite{kapoor2023leakage}。第一，在查看 held-out DOI 标签之前，必须冻结评分、候选生成、前驱体库、路线族阈值、划分规则和提交的排名表。第二，每个 held-out 行必须连接到 \texttt{frozen\_route\_rankings.csv} 中归档的提交排名，其中 \texttt{matched\_route\_id} 映射到冻结排名的 \texttt{route\_id}，且排名提交时间早于标签释放。第三，文献路线、路线族、best-top-$k$、预期标签、人工标签和裁决标签都必须排除在评分输入之外。在填充后的冻结排名和 held-out 路线表通过这些检查之前，blind-style recovery 仍是验证升级，而不是数值证据。

\subsection*{\invdesroute{} scoring, ablation and diversity-aware display}
% 中文：\invdesroute{} 评分、消融和 diversity-aware 展示

The \invdesroute{} score is a route-recommendation score applied after reaction enumeration and energy evaluation. It does not modify PET-OAM-XL energies. The score combines PET-form-aligned reaction favorability, precursor priors, route-template priors, practical-route status, coefficient complexity, phase availability and condition simplicity:
% 中文：\invdesroute{} score 是在反应枚举和能量评估之后应用的路线推荐分数。它不修改 PET-OAM-XL 能量。该分数结合 PET-form-aligned 反应有利性、前驱体先验、路线模板先验、practical-route 状态、系数复杂度、相可用性和条件简洁性：
\begin{equation}
S_{\mathrm{IDR}}
=0.18S_E
+0.25S_P
+0.25S_T
+0.12S_R
+0.08S_C
+0.07S_A
+0.05S_Q
+B_T ,
\end{equation}
where $S_E$, $S_P$, $S_T$, $S_R$, $S_C$, $S_A$ and $S_Q$ denote energy favorability, precursor prior, route template, practical-route status, complexity, phase availability and condition simplicity. A template bonus $B_T$ of 0.06 is added for near-canonical templates ($S_T\geq0.95$) and 0.03 for strong templates ($S_T\geq0.90$), with the final score clipped to [0,1]. The binned energy term assigns $S_E=1.00$, 0.88, 0.64, 0.42, 0.20 and 0.35 for reactions that are strongly exergonic ($\Delta E\leq-1.0$ eV), weakly exergonic, near thermoneutral, uphill, strongly uphill or energy-unavailable, respectively; the unavailable case also carries a risk flag.
% 中文：其中 $S_E$、$S_P$、$S_T$、$S_R$、$S_C$、$S_A$ 和 $S_Q$ 分别表示能量有利性、前驱体先验、路线模板、practical-route 状态、复杂度、相可用性和条件简洁性。近规范模板（$S_T\geq0.95$）加入 0.06 的模板加分 $B_T$，强模板（$S_T\geq0.90$）加入 0.03，最终分数截断到 [0,1]。分箱能量项对强放热（$\Delta E\leq-1.0$ eV）、弱放热、近热中性、上坡、强上坡和能量不可用的反应分别赋予 $S_E=1.00$、0.88、0.64、0.42、0.20 和 0.35；能量不可用的情况还带有风险标记。

The non-energy components encode route chemistry: precursor families, route templates, coefficient and formula complexity, phase availability, gas reservoirs, gas byproducts, atmosphere-control requirements and likely precursor decomposition. They alter display priority while preserving thermodynamic records and risk flags.
% 中文：非能量组件编码路线化学，包括前驱体家族、路线模板、系数和化学式复杂度、相可用性、气体库、气体副产物、气氛控制需求和可能的前驱体分解。它们改变展示优先级，同时保留热力学记录和风险标记。

Ranking ablations used the 59-route overlap set, 86-route literature-seeded extension and combined 145-route comparison. PET-energy-only, practical-first, full \invdesroute{} score and diversity-aware display modes were evaluated by the best route-family similarity within top-1, top-5, top-10 and top-20 windows. Workflow-level baselines are treated as a separate validation upgrade because target-stability-only and route-ignorant policies rank targets, whereas energy-only and route-aware policies rank routes. The required comparison freezes four policies, the shared 145-route evaluation set, the candidate set, the top-$k$ window, the route-exposure source, score-input allow/deny lists and pre-label ranking flags. It must report both policy-unit rank and route-exposure rank, yielding 580 included records for the 145-route top-5 design. Until those records are populated, the present result is limited to route-level visibility.
% 中文：排序消融使用 59 条重叠路线、86 条文献种子扩展路线以及合并后的 145 条路线比较。PET-energy-only、practical-first、完整 \invdesroute{} score 和 diversity-aware display 模式通过 top-1、top-5、top-10 和 top-20 窗口内的最佳路线族相似性进行评价。工作流级基线被视为单独的验证升级，因为 target-stability-only 和 route-ignorant 策略排序的是目标，而 energy-only 和 route-aware 策略排序的是路线。所需比较会冻结四种策略、共享的 145 条路线评价集、候选集、top-$k$ 窗口、路线暴露来源、评分输入 allow/deny 列表和标签前排名标记。它必须同时报告 policy-unit rank 和 route-exposure rank，在 145 条路线的 top-5 设计中得到 580 条纳入记录。在这些记录被填充之前，当前结果仅限于路线级可见性。

For the literature-seeded extension, additional targets were selected from the same DOI-backed route pool when they were absent from the 120-target benchmark, passed route filters, had cached MP phase tables and had reusable PET-OAM-XL energies. The 53 targets were analyzed separately to test portability to experimentally occurring compositions, not to estimate the full distribution of inorganic synthesis routes.
% 中文：对于文献种子扩展，额外目标从同一带 DOI 支撑路线池中选择，条件是它们不在 120 目标 benchmark 中、通过路线过滤、具有缓存 MP 相表，并具有可复用的 PET-OAM-XL 能量。这 53 个目标被单独分析，用于测试向实验中出现组成的可迁移性，而不是估计无机合成路线的完整分布。

\subsection*{Finite-temperature gas-reservoir sensitivity}
% 中文：有限温气体库敏感性

Finite-temperature effects were treated through first-order gas-reservoir sensitivity. Routes containing CO$_2$, O$_2$, NH$_3$ or H$_2$ were corrected with ideal-gas chemical-potential terms at 700, 1000 and 1300 K using representative partial pressures. Corrections were applied only to gas species; solid vibrational, phase-transition, defect, surface, pressure--volume and kinetic terms were not included.
% 中文：有限温效应通过一阶气体库敏感性处理。含 CO$_2$、O$_2$、NH$_3$ 或 H$_2$ 的路线使用代表性分压，在 700、1000 和 1300 K 下加入理想气体化学势项进行校正。校正只应用于气体物种；不包括固体振动、相变、缺陷、表面、压力--体积和动力学项。

\subsection*{Claim-to-evidence release checks}
% 中文：主张到证据的发布检查

The release package implements a claim-to-evidence record rather than an additional result. Each major claim and captioned panel links to source tables, statistic scripts, check gates, allowed interpretation and boundary not claimed, following expectations for reproducible computational materials research, machine-learning reproducibility and dataset documentation \cite{wilkinson2016fair,lejaeghere2016dft_reproducibility,coudert2017reproducible_research,pineau2021reproducibility,gebru2021datasheets}. Template-mode gates check schemas, provenance rows and claim boundaries during manuscript preparation. Submission-mode gates must be run on the archived release and require populated source tables, frozen configurations, non-placeholder identifiers, regenerated figure-source records and recomputed manuscript statistics from one release directory.
% 中文：发布包实现的是主张到证据的记录，而不是额外结果。每个主要主张和带图注的 panel 都连接到源表、统计脚本、检查门、允许解释和未主张边界，并遵循可复现计算材料研究、机器学习可复现性和数据集文档的要求 \cite{wilkinson2016fair,lejaeghere2016dft_reproducibility,coudert2017reproducible_research,pineau2021reproducibility,gebru2021datasheets}。Template-mode gates 在手稿准备期间检查 schema、provenance rows 和 claim boundaries。Submission-mode gates 必须在归档发布版本上运行，并要求填充后的源表、冻结配置、非占位标识符、重新生成的 figure-source records，以及从同一 release directory 重新计算的手稿统计量。



\section*{Data availability}
% 中文：数据可用性

The Materials Project structures and thermodynamic metadata used by \invdesroute{} for phase retrieval are available from Materials Project subject to its access terms \cite{jain2013materialsproject}. The DOI-backed literature-route validation set was derived from the public text-mined solid-state synthesis dataset of Kononova \textit{et al.} \cite{kononova2019textmined,cedergrouphub_textmined_synthesis_public} after the filtering criteria described in Methods. The processed data generated in this study, including model-evidence benchmark tables, the 120-target route-network benchmark, PET/MP-form comparable reaction tables, DOI-backed literature-route audit tables, route-family similarity scores, recommendation-ablation results, precursor-class diagnostics, confidence-label records, refinement queues, validation summaries, reproducibility manifests and figure-source tables, are available in the InvDesRoute data archive at \url{ZENODO_DOI_OR_URL}. Materials Project API keys, raw Materials Project database content, full CIF caches and model checkpoints are not redistributed in the archive; users should obtain these resources from their original providers under the applicable access terms.
% 中文：\invdesroute{} 用于相检索的 Materials Project 结构和热力学元数据可在遵守 Materials Project 访问条款的前提下获得 \cite{jain2013materialsproject}。带 DOI 支撑的文献路线验证集来自 Kononova \textit{et al.} 的公开文本挖掘固相合成数据集 \cite{kononova2019textmined,cedergrouphub_textmined_synthesis_public}，并经过 Methods 中描述的过滤标准处理。本文生成的处理后数据，包括模型证据 benchmark 表、120-target 路线网络 benchmark、PET/MP-form 可比反应表、DOI 支撑的文献路线审计表、路线族相似度分数、推荐消融结果、前驱体类别诊断、置信标签记录、精修队列、验证汇总、可复现性清单和作图源数据，均存放在 InvDesRoute 数据归档中：\url{ZENODO_DOI_OR_URL}。归档不重新分发 Materials Project API key、Materials Project 原始数据库内容、完整 CIF cache 或模型 checkpoint；用户应按照原始提供方的适用访问条款获取这些资源。

\section*{Code availability}
% 中文：代码可用性

The \invdesroute{} code for phase retrieval, PET energy caching, balanced reaction enumeration, MP-form reference construction, practical-route filtering, route scoring, literature-route comparison, diversity-aware recommendation, precursor-class diagnostics and gas-reservoir sensitivity analysis is available at \url{GITHUB_OR_ARCHIVE_URL}. The archived code snapshot corresponding to this manuscript release is available at \url{CODE_ARCHIVE_DOI_OR_URL}. The repository contains workflow scripts, example formula and CIF inputs, environment specifications, frozen configuration notes, reproducibility manifests and commands for regenerating the benchmark tables, supplementary tables and figure-source data. Materials Project queries require user-provided Materials Project credentials. PET-OAM-XL and other model checkpoints are not redistributed and should be obtained from their original sources.
% 中文：用于相检索、PET 能量缓存、配平反应枚举、MP-form 参考构建、practical-route 过滤、路线评分、文献路线比较、diversity-aware 推荐、前驱体类别诊断和气体库敏感性分析的 \invdesroute{} 代码可在 \url{GITHUB_OR_ARCHIVE_URL} 获取。与本文发布版本对应的归档代码快照可在 \url{CODE_ARCHIVE_DOI_OR_URL} 获取。该仓库包含工作流脚本、示例化学式和 CIF 输入、环境规格、冻结配置说明、可复现性清单，以及用于再生成 benchmark 表、补充表和作图源数据的命令。Materials Project 查询需要用户自行提供 Materials Project 凭据。PET-OAM-XL 和其他模型 checkpoint 不随本文重新分发，应从其原始来源获取。




\vspace{6pt}
\noindent\textbf{Acknowledgement:} This work is supported by the Beijing Natural Science Foundation (No. Z250005), the National Natural Science Foundation of China (Nos. 62476278 and 11934020), and the National Key R\&D Program of China (No. 2024YFA1408601). Computational resources were provided by the Physical Laboratory of High Performance Computing at Renmin University of China. We thank Zhenfeng Ouyang and Jing Jiang for helpful discussions.
% 中文：致谢：本工作得到北京市自然科学基金（No. Z250005）、国家自然科学基金（Nos. 62476278 and 11934020）和国家重点研发计划（No. 2024YFA1408601）的支持。计算资源由中国人民大学高性能计算物理实验室提供。我们感谢 Zhenfeng Ouyang 和 Jing Jiang 的有益讨论。

\noindent\textbf{Corresponding authors:} Correspondence and requests for materials should be addressed to Ze-Feng Gao (zfgao@ruc.edu.cn) and Zhong-Yi Lu (zlu@ruc.edu.cn).
% 中文：通讯作者：通信和材料请求应发送给 Ze-Feng Gao (zfgao@ruc.edu.cn) 和 Zhong-Yi Lu (zlu@ruc.edu.cn)。

\noindent\textbf{Competing interests:} The authors declare no competing interests. 

\noindent\textbf{Supplementary materials:} Supplementary materials are provided separately.

\clearpage


\bibliographystyle{unsrtnat}
\bibliography{references.bib}

% add extended data figures/tables below





\end{document}
